\section{The Four-Module Decomposition}
\label{sec:design-modules}
% Authoring brief for this section lives in section.md (§5.3).

The architectural style of \cref{sec:design-style} fixes the direction of dependencies but not the concerns they organise.
This section opens the core and identifies those concerns.
We decompose the platform's adaptive behaviour into four parts (sensing, eye-movement analysis, decision, and intervention), chosen not by technical tier but along the axis of replaceability that the dominant driver D1 demands.
Each part is a module that hides its internals behind a contract, and the four contracts together are the structural claim that RQ1 puts to the test.

The four concerns form a pipeline that follows the data, from raw signal to committed change.
Sensing is the source: it produces a continuous stream of gaze samples, and its contract is defined over those samples rather than over any device, so that the Tobii tracker and a simulated, mouse-cursor-driven source are interchangeable behind it (D5).
Analysis consumes that stream and projects it into oculomotor events (the fixations, saccades, and regressions that constitute reading behaviour), abstracting the particular identification algorithm behind its contract \cite{salvucci2000identifying}.
Decision consumes the analysed state together with the session context and returns, at most, a proposal to intervene; this is the seam at which an external, AI-driven provider stands in for the built-in rule-based strategy, and it is the structural origin of the out-of-scope constraint C3.
Intervention is the sink: it validates a proposal and then commits a context-preserving change to the presented text at a controlled boundary, which is the property that RQ3 and driver~D7 put to the test.
\Cref{fig:design-modules} shows the four as ports on a common core, a structure we develop in the remainder of this section.

\begin{figure}[htbp]
\centering
\resizebox{0.92\linewidth}{!}{%
\begin{tikzpicture}[
  hexcore/.style={draw=navyblue, very thick, fill=blue!7},
  hexdom/.style={draw=navyblue, thick, fill=blue!16},
  port/.style={rectangle, rounded corners=2pt, draw=dtured, very thick, fill=dtured!12, align=center, text width=1.7cm, minimum height=0.8cm, inner sep=2pt, font=\small},
  impl/.style={rectangle, rounded corners=2pt, draw=black!55, thick, fill=grey!30, align=center, text width=2.4cm, minimum height=0.7cm, inner sep=2pt, font=\scriptsize},
  inlink/.style={->, thick, black!60},
  cyc/.style={->, semithick, navyblue!75, shorten >=1pt, shorten <=1pt},
  note/.style={font=\scriptsize\itshape, text=black!70, align=center},
]
% outer hexagon: the application core boundary (ports sit on it)
\draw[hexcore] (-1.5,2.6) -- (1.5,2.6) -- (3.0,0) -- (1.5,-2.6) -- (-1.5,-2.6) -- (-3.0,0) -- cycle;
% inner hexagon: the domain
\draw[hexdom] (-0.65,1.13) -- (0.65,1.13) -- (1.3,0) -- (0.65,-1.13) -- (-0.65,-1.13) -- (-1.3,0) -- cycle;
\node[font=\small\bfseries, align=center] at (0,0) {Application\\core};
\node[note] at (0,1.82) {domain + reading runtime};
% ports, sitting on the hexagon boundary
\node[port] (psen) at (-3.0, 0)  {{\bfseries Sensing}\\{\scriptsize port (1)}};
\node[port] (pana) at (0, 2.6)   {{\bfseries Analysis}\\{\scriptsize port (2)}};
\node[port] (pdec) at (3.0, 0)   {{\bfseries Decision}\\{\scriptsize port (3)}};
\node[port] (pint) at (0, -2.6)  {{\bfseries Intervention}\\{\scriptsize port (4)}};
% adapters and providers, outside the core
\node[impl] (tobii)  at (-6.3, 0.6)  {Tobii adapter};
\node[impl] (sims)   at (-6.3,-0.6)  {Simulated (mouse) source};
\node[impl] (anab)   at (-1.8, 4.4) {Built-in strategy};
\node[impl] (anae)   at ( 1.8, 4.4) {External analysis provider};
\node[impl] (decr)   at (6.3, 1.0)  {Rule-based strategy};
\node[impl] (dece)   at (6.3,-1.0)  {External decision provider};
\node[impl] (intb)   at (-1.8,-4.4) {Built-in modules};
\node[impl] (inte)   at ( 1.8,-4.4) {Additional modules};
% inward dependency arrows (adapters depend on the ports)
\draw[inlink] (tobii.east)  -- (psen.west);
\draw[inlink] (sims.east)   -- (psen.west);
\draw[inlink] (anab.south)  -- (pana.north);
\draw[inlink] (anae.south)  -- (pana.north);
\draw[inlink] (decr.west)   -- (pdec.east);
\draw[inlink] (dece.west)   -- (pdec.east);
\draw[inlink] (intb.north)  -- (pint.south);
\draw[inlink] (inte.north)  -- (pint.south);
% adaptive-cycle order between the ports (sensing -> analysis -> decision -> intervention -> ...)
\draw[cyc] (psen.north) to[bend right=12] (pana.west);
\draw[cyc] (pana.east)  to[bend right=12] (pdec.north);
\draw[cyc] (pdec.south) to[bend right=12] (pint.east);
\draw[cyc] (pint.west)  to[bend right=12] (psen.south);
% legend: the two arrow kinds
\node[note, align=left, anchor=west] at (-6.7,-3.4) {Grey arrows: dependency direction\\(inward, to the core that owns the ports).\\Navy arrows: the adaptive-cycle order.};
\end{tikzpicture}%
}
\caption{The four-module decomposition. Each concern (sensing, analysis, decision, intervention) is a port owned by the application core, while concrete implementations sit outside the core and depend inward upon the ports (solid arrows), so that the core depends on none of them. The numbers and the curved navy arrows between the ports trace the data-flow order of one adaptive cycle, from sensing (1) through analysis (2) and decision (3) to intervention (4) and back to sensing. Each port resolves its implementation through a registry, so that adding an implementation is an additive change.}
\label{fig:design-modules}
\end{figure}

The three concerns that carry intelligence (analysis, decision, and intervention) deliberately share one contract shape.
Each is identified by a stable provider identity and exposes a single operation that receives an immutable snapshot of the relevant session state and returns an optional result.
The snapshot is the load-bearing choice: a module is handed a copy of what it needs to see, not a reference to the live session, so that it can neither reach into nor mutate the core's state.
This is what makes a module replaceable without coordination and testable in isolation, since its entire interaction with the system is one transformation from a snapshot to a result.
The uniformity across the three seams is itself a design decision rather than a coincidence, because a contributor who has implemented one kind of module already understands the shape of the others.

These contracts are owned by the core, and that ownership is what gives the dependency rule its force.
The interfaces belong to the application core, whereas the concrete implementations (the Tobii adapter, the rule-based and external strategies, and the individual intervention modules) live outside it and depend inward upon them \cite{martin1996dip}.
The core in turn depends on none of these implementations, being expressed entirely in terms of its own ports and domain types.
Read concentrically, as in \cref{sec:design-style}, this places the domain and application at the centre and every replaceable detail at the rim, with all dependencies pointing inward.
``Independently replaceable'' can now be stated precisely: any implementation may be exchanged for another that satisfies the same port without a single change to the core, which is exactly the property evaluated against RQ1.

Replaceability becomes extensibility through one further element at each seam, a registry that resolves an implementation by its provider identity.
Because resolution is by identity rather than by type, adding a new analysis algorithm, decision provider, or intervention is an additive change: a new implementation of an existing port, registered alongside the others, with no edit to the reading runtime that drives them.
The platform ships a built-in implementation at each seam and, for analysis and decision, a second implementation that forwards the same contract to an out-of-process provider, so that the extension path is not hypothetical but already exercised by the system itself.
This mechanism, a strategy selected at runtime from a registry, is how the modularity named in the title of this thesis is delivered \cite{gamma1994designpatterns}.

One consequence of the same discipline deserves emphasis, because it is what makes the architecture practical to develop and to defend.
Since the sensing port is written in terms of gaze samples and not Tobii types, the analysis, decision, and intervention layers above it run unchanged when the source is the simulated, mouse-cursor stream rather than the Tobii device.
The platform therefore does not require the eye-tracking hardware in order to exercise the adaptive pipeline from end to end, which supports development, automated testing, and demonstration away from the device (D5), and which allows each layer to be validated in isolation.

The decomposition described here is static: it names the four concerns, their contracts, and the direction of their dependencies.
What it does not yet show is how the four behave together under the timing constraints of a live session.
\Cref{sec:design-loop} sets these ports in motion as the real-time adaptive loop.
